% \usetikzlibrary{arrows.meta}

% -------------------------------------------------
% Styly
% -------------------------------------------------
\tikzstyle{faactive}=[fill=red!20]

\tikzstyle{compnode}=[%
    draw,
    rounded corners=2pt,
    thick,
    inner xsep=2mm,
    inner ysep=1mm,
    font=\small
]

\tikzstyle{compactive}=[%
    compnode,
    fill=red!20
]

\tikzstyle{compedge}=[thick]

\tikzstyle{resultbox}=[
    draw,
    rounded corners=2pt,
    thick,
    fill=green!20,
    inner xsep=3mm,
    inner ysep=2mm,
    font=\small
]

\newcommand{\FAResult}[2][]{%
    \node[resultbox,#1] at (7.8,-1.2) {#2};
}

% -------------------------------------------------
% Pomocná makra pro zobrazení vstupního slova
% -------------------------------------------------

% #1 ... přečtená část
% #2 ... aktuálně čtený znak
% #3 ... zbytek slova
\newcommand{\FAInput}[3]{%
    \node[font=\large] at (2.4,1.45) {$#1\textcolor{red}{#2}#3$};
}

% Varianta po dočtení celého slova
% #1 ... celé přečtené slovo
\newcommand{\FAInputFinished}[1]{%
    \node[font=\large] at (2.4,1.45) {$#1\;\textcolor{red}{\lambda}$};
}

% -------------------------------------------------
% Uzly a hrany stromu výpočtu
% -------------------------------------------------

% #1 ... volitelné tikz volby
% #2 ... jméno uzlu
% #3 ... souřadnice
% #4 ... obsah uzlu
\newcommand{\CompNode}[4][]{%
    \node[compnode,#1] (#2) at #3 {$#4$};
}

% #1 ... volitelné tikz volby
% #2 ... odkud
% #3 ... popisek hrany
% #4 ... kam
\newcommand{\CompEdge}[4][]{%
    \draw[compedge,#1] (#2) -- node[midway,font=\scriptsize,fill=white,inner sep=1pt] {$#3$} (#4);
}

% -------------------------------------------------
% Makro pro jeden slide animace
% -------------------------------------------------
% #1 ... styl stavu q0
% #2 ... styl stavu q1
% #3 ... styl stavu q2
% #4 ... kód zobrazující vstupní slovo
% #5 ... kód stromu výpočtu
\newcommand{\AutomatonComputationFrame}[5]{%
\begin{frame}{Výpočet deterministického automatu}{Příklad}
% \centering
\begin{tikzpicture}[>=stealth]

    % Vstupní slovo
    #4

    % Automat
    \initialstate[#1]{q0}{(0,0)}{$q_0$}{(-1.2,0)}{west}
    \state[#2]{q1}{(2.4,0)}{$q_1$}
    \acceptingstate[#3]{q2}{(4.8,0)}{$q_2$}

    \transitionarrow{q0}{above}{$0$}{q1}
    \transitionarrow{q1}{above}{$1$}{q2}
    \transitionarrow[loop above]{q2}{above}{$0,1$}{q2}

    % Popisek stromu
    \node[anchor=west,font=\normalsize] at (-0.4,-1.15) {Strom výpočtu:};

    % Strom výpočtu
    #5

\end{tikzpicture}
\end{frame}
}